\documentclass[11pt,a4paper]{article}
\usepackage[T1]{fontenc}
\usepackage[utf8]{inputenc}
\usepackage[english]{babel}
\usepackage[margin=27mm]{geometry}
\usepackage{amsmath,amssymb,amsthm,mathtools,bm}
\usepackage{booktabs,longtable,array}
\usepackage{xcolor,listings,microtype}
\usepackage[hidelinks]{hyperref}
\usepackage{enumitem}
\setlist{nosep}
\setlength{\emergencystretch}{2em}
\newcommand{\R}{\mathbb R}
\newcommand{\E}{\mathbb E}
\newcommand{\tr}{\operatorname{tr}}
\newcommand{\rank}{\operatorname{rank}}
\newcommand{\diag}{\operatorname{diag}}
\newcommand{\col}{\operatorname{col}}
\newcommand{\Gr}{\operatorname{Gr}}
\newcommand{\norm}[1]{\left\lVert#1\right\rVert}
\newcommand{\code}[1]{\texttt{\detokenize{#1}}}
\newtheorem{proposition}{Proposition}
\newtheorem{remark}{Remark}
\lstset{basicstyle=\ttfamily\small,columns=fullflexible,keepspaces=true,
  breaklines=true,frame=single,rulecolor=\color{black!20},showstringspaces=false}
\title{The Current Average Derivative Procedure\\
for Multi-Index Dimension Reduction\\[3pt]
\large Complete Algorithm and the \texttt{grassman.py} Solver}
\author{Implementation-based mathematical description}
\date{Source snapshot: 23 October 2026}
\begin{document}
\maketitle
\begin{abstract}
This document describes the complete global multi-index ADP pipeline in the
current repository when its inner solver is explicitly set to
\code{ADP/solver/grassman.py::solve}. It covers the statistical motivation,
centers, local-gradient initialization, adaptive kernel geometry, random
direction sketches, centered local moments, variable projection on the
Grassmann manifold, numerical safeguards, spectrum updates, bandwidth
continuation, and selection of the returned estimate. The main path is the
normalized \code{estimator="new"} configuration with the current twenty-step
Grassmann solver. Mathematical identities are distinguished from finite-budget
optimization choices and from statistical assumptions. The description follows
the live code; it does not assert global convergence or a recovery theorem.
\end{abstract}
\tableofcontents
\newpage

\section{Statistical target and the two levels of the algorithm}
\label{sec:target}
We observe $(x_i,y_i)_{i=1}^n$, with $x_i\in\R^d$ and $y_i\in\R$.
The motivating model is
\begin{equation}\label{eq:model}
 y_i=f(x_i)+\varepsilon_i,\qquad
 \E[\varepsilon_i\mid x_i]=0,\qquad
 f(x)=\psi(B_\star^\top x),\qquad B_\star^\top B_\star=I_m,
\end{equation}
where the index dimension $1\le m<d$ is supplied by the user. The target is
the effective dimension reduction (EDR) subspace $\col(B_\star)$, not an
oriented basis and not the link function $\psi$.

If $\psi$ is differentiable, then
\begin{equation}\label{eq:gradient-model}
 \nabla f(x)=B_\star\nabla\psi(B_\star^\top x).
\end{equation}
Thus local gradients belong to a common $m$-dimensional space. If their
coefficient second moment is positive definite, their span identifies that
space. Smoothness, sufficiently informative neighborhoods, excitation of all
$m$ directions, and adequate control of noise are statistical requirements;
they are not consequences of a small numerical solver residual.

ADP alternates two different operations:
\begin{enumerate}
\item An \emph{outer statistical update} uses the current EDR estimate to
choose anisotropic neighborhoods, draws a finite direction sketch, and builds
new local moment equations.
\item An \emph{inner numerical update} freezes those moments and approximately
minimizes a profiled least-squares objective over an EDR subspace.
\end{enumerate}
The outer update changes the objective itself. Decrease of one frozen inner
objective does not imply decrease of a common objective across outer steps,
nor improvement in distance to $\col(B_\star)$.

\paragraph{Scope and actual entry point.}
This is the shared global multi-index pipeline, not the separate manifold
model with one subspace per center. The public multi-index class uses LSMR
unless a solver is supplied. The algorithm described here is selected by
\begin{lstlisting}[language=Python]
from ADP import ADP_Config, ADP_multi_index, ADP_solver
from ADP.solver.grassman import solve

cfg = ADP_Config()  # initializer ridge and outer parameters
model = ADP_multi_index(m, cfg, ADP_solver(solve)).fit(X, y)
B_hat = model.basis_       # shape (d, m), orthonormal columns
Z = model.transform(X)     # X @ B_hat
\end{lstlisting}
The data matrix \code{X} stores observations in rows. The fitted object
provides a projection; fitting a downstream regression function on $Z$ is
a separate task.

\section{Notation, configuration, and data protocol}\label{sec:notation}
\begin{center}
\begin{tabular}{lll}
\toprule
Symbol & Shape & Meaning\\
\midrule
$X$, $y$ & $n\times d$, $n$ & observations and responses\\
$c_j$, $j=1,\ldots,J$ & $d$ & localization centers\\
$B_k$ & $d\times m$ & outer basis, $B_k^\top B_k=I_m$\\
$\ell_k$ & $m$ & normalized localization spectrum\\
$h_k$, $\alpha_k$ & scalars & bandwidth and anisotropy\\
$\Phi_j$ & $p\times d$ & rows $\phi_{js}^\top$ of unit directions\\
$N_j$, $\mu_j$, $\bar y_j$ & scalar, $d$, scalar & local mass and means\\
$I_j$, $U_j$ & $p$, $p\times d$ & directional response and operator\\
$Y$, $g_j$ & $d\times m$, $m$ & inner basis and local coefficients\\
\bottomrule
\end{tabular}
\end{center}
The code passes a row basis $P=B^\top$ of shape $(m,d)$ to a solver and
returns the same layout. This document uses column bases $B$ and $Y$ to
make the geometry explicit. The array called \code{I} in the solver stores
the vectors $I_j$; $I_m$ in a formula denotes the identity matrix.

The default values relevant to the main path are listed below. Derived
sizes are evaluated after the data dimension is known.
\begin{center}
\begin{tabular}{ll}
\toprule
Outer setting & Current value\\
\midrule
\code{seed} & $42$\\
\code{N_loc}, \code{N_lin} & $10$, default $2d$\\
\code{N_J}, \code{N_phi} & default $n$, default $\min(N_{\rm loc},d)$\\
\code{index_init}, \code{estimator} & \code{local}, \code{new}\\
\code{direction_mode}, \code{redraw_directions} & \code{auto}, \code{True}\\
\code{multi_tensor}, \code{select_step} & \code{orthogonal}, \code{best}\\
\code{training_set}, \code{center_displacement} & \code{all}, $0$\\
\code{a}, \code{h_min}, \code{outer_steps} & $\sqrt{2}$, $1$, no explicit cap\\
\code{local_ridge}, \code{lambda_penalty} & $10^{-8}$, $0.05$\\
\bottomrule
\end{tabular}
\end{center}
Multi-index fitting requires $p>m$. Non-random initialization additionally
requires $n>d+1$; sample sizes, requested masses, and center counts are
validated. An average-mass constraint does not ensure every center has
positive mass: an empty local neighborhood is rejected by the statistics code.

\paragraph{Centers and random streams.}
Centers are sampled from observation indices without replacement; if $J=n$,
all observations are used. Three child streams of \code{SeedSequence(seed)}
separate center selection, initialization, and direction generation. Random
initialization uses its stream; the optional pilot initializer uses the
configuration seed. Optional displaced centers have the form
$c_j=x_{i_j}+\nu\sigma_X z_j$, where $z_j\sim N(0,I_d)$ and the default
$\sigma_X$ is the root mean feature variance. Center response labels remain
$y_{i_j}$ even when $c_j$ is displaced.

The default \code{training_set} is \code{all}: every observation contributes.
With the \code{exclude_centers} option, the original sampled indices are removed from
the training observations, while their responses are retained for outer-step
selection. The default uses the same observations for training and selection;
it is not an independent validation split. Below, sums over $i$ refer to the
resulting training set, of size $n_{\rm tr}$.

\section{Initialization by local gradient estimation}\label{sec:init}
The default initializer first chooses an isotropic bandwidth $h_{\rm lin}$
with average kernel mass at least $N_{\rm lin}$. At each center, define
$w^{\rm lin}_{ji}=K(\norm{x_i-c_j}^2/h_{\rm lin}^2)$ and solve
\begin{equation}\label{eq:init-fit}
 (\hat a_j,\hat b_j)\in\arg\min_{a\in\R,b\in\R^d}
 \sum_i w^{\rm lin}_{ji}\bigl[y_i-a-b^\top(x_i-c_j)\bigr]^2
 +\lambda_{\rm init}\norm{b}^2.
\end{equation}
The intercept is not penalized. The code uses centered SVD or augmented
least squares, with numerical rank checks; it does not require an inverse of
a local covariance. Here $\lambda_{\rm init}$ is
\code{ADP_Config.local_ridge}, whose default is $10^{-8}$.

Let $L_{\rm init}$ have rows
$\sqrt{N_j^{\rm lin}}\hat b_j^\top$ for \code{estimator="new"}.
The first $m$ right singular vectors of $L_{\rm init}$ give $B_0$. This is
equivalent to the leading eigenspace of
\begin{equation}
 \sum_j N_j^{\rm lin}\hat b_j\hat b_j^\top,
\end{equation}
without forming that $d\times d$ matrix. Insufficient numerical rank is
reported as failure to identify the requested index dimension.

The \code{local-cv} initializer selects a ridge by weighted leave-one-out
error before the gradient PCA; \code{random} uses a Gaussian QR basis;
\code{pilot} uses an orthonormalized shallow MLP first-layer space.
These are explicit alternative initializers. In all cases, the outer loop
starts with $\ell_0=(1,\ldots,1)$ and $\alpha_0=1$. The initializer's
gradient spectrum is recorded as metadata but does not replace $\ell_0$.

\section{Adaptive localization and bandwidth continuation}\label{sec:local}
\subsection{The implemented kernel and localization geometry}
At the start of outer step $k$, for $\delta_{ji}=x_i-c_j$, put
\begin{equation}\label{eq:components}
 z_{ji}=B_k^\top\delta_{ji},\qquad
 r_{ji}^2=\norm{\delta_{ji}}^2-\norm{z_{ji}}^2,\qquad
 s_{ji}^2=\sum_{a=1}^m\ell_{k,a}z_{ji,a}^2.
\end{equation}
In floating point, negative roundoff in $r_{ji}^2$ is clipped to zero.
The two supported tensors give
\begin{align}
 t_{ji}^{\perp}&=\frac{\alpha_k^2 r_{ji}^2+s_{ji}^2}{h_k^2},
 &A_k^{\perp}&=\frac{\alpha_k^2(I_d-B_kB_k^\top)
                  +B_k\diag(\ell_k)B_k^\top}{h_k^2},\label{eq:orthogonal}\\
 t_{ji}^{\rm full}&=\frac{\alpha_k^2\norm{\delta_{ji}}^2+s_{ji}^2}{h_k^2},
 &A_k^{\rm full}&=\frac{\alpha_k^2I_d
                  +B_k\diag(\ell_k)B_k^\top}{h_k^2}.
 \label{eq:full}
\end{align}
Thus $t_{ji}=\delta_{ji}^\top A_k\delta_{ji}$ and $w_{ji}=K(t_{ji})$.
These matrices are mathematical notation; the implementation evaluates the
quadratic forms through projections rather than constructing dense tensors.

The exact default kernel convention is
\begin{equation}\label{eq:kernel}
 K(t)=\max(1-t^2,0),\qquad t\ge0.
\end{equation}
Its input here is already a \emph{squared} normalized distance. In an
isotropic neighborhood, the weight is therefore
$\max(1-\norm{\delta_{ji}}^4/h_k^4,0)$, not
$\max(1-\norm{\delta_{ji}}^2/h_k^2,0)$. This convention follows
\code{epanechnikov} and its actual callers.

The first multi-index step always uses the orthogonal tensor, including
when the requested \code{multi_tensor} is \code{full}. Since
$\alpha_0=1$ and $\ell_0=\bm 1$, that first neighborhood is isotropic.
Subsequent steps use the configured tensor. The default orthogonal tensor
broadens neighborhoods in the estimated irrelevant directions as $\alpha$
decreases. Small $\ell_a$ also broadens a neighborhood along its $a$th
estimated principal direction.

\subsection{Choosing \texorpdfstring{$h$ and $\alpha$}{h and alpha}}
The initial outer bandwidth $h_0\ge h_{\min}$ is the smallest feasible
isotropic bandwidth, up to bisection accuracy, for
\begin{equation}\label{eq:h-search}
 \frac1J\sum_{j,i}K\left(\frac{\norm{x_i-c_j}^2}{h_0^2}\right)
 \ge N_{\rm loc}.
\end{equation}
The search doubles an upper bracket and then bisects. Local-gradient
initialization uses the same type of search, but targets $N_{\rm lin}$
and uses a small positive lower bound rather than $h_{\min}$.

After estimating $B_{k+1}$ and $\ell_{k+1}$, the proposed next bandwidth is
$h_{k+1}=h_k/a$, with $a>1$. With that new geometry held fixed, the code
selects the largest feasible $\alpha\in[0,1]$ satisfying
\begin{equation}\label{eq:alpha-search}
 \frac1J\sum_{j,i}w_{ji}(B_{k+1},\ell_{k+1},h_{k+1},\alpha)
 \ge N_{\rm loc}.
\end{equation}
For the default kernel, mass is non-increasing in $\alpha$:
if $\alpha=1$ is feasible it is returned; if $\alpha=0$ is infeasible,
continuation stops; otherwise bisection returns the feasible lower endpoint
to tolerance $\sqrt{\epsilon_{\rm mach}}$. There is no constraint that the
new $\alpha$ be below the previous one. Custom kernels must support the
monotonicity assumed by these searches.

Compact-support screening removes only pairs proved to have zero weight
throughout the remaining search interval. It preserves the chosen
finite-sample estimator. The target is average \emph{kernel mass}, not
neighbor count and not effective sample size.

\section{Random directions and centered moment equations}\label{sec:moments}
\subsection{The finite direction sketch}
For \code{estimator="new"} in multi-index mode,
\code{direction_mode="auto"} resolves to \code{isotropic}:
\begin{equation}\label{eq:isotropic}
 \phi_{js}=z_{js}/\norm{z_{js}},\qquad z_{js}\sim N(0,I_d),
 \qquad s=1,\ldots,p.
\end{equation}
Directions are independently redrawn at each outer step by default.
Setting \code{redraw_directions} to \code{False} reuses the first sketch, so the
random finite-sample procedure changes.

The explicit \code{localized} option instead draws independent
$z\sim N(0,I_d)$ and $\xi\sim N(0,I_m)$ and normalizes
\begin{equation}\label{eq:localized-directions}
 v=\alpha_k(I_d-B_kB_k^\top)z
      +B_k\diag(\sqrt{\ell_k})\xi,\qquad \phi=v/\norm{v}.
\end{equation}
This option changes the sketch distribution; it is not the default
normalized multi-index path.

\subsection{Weighted centering and the exact empirical quantities}
For each center define
\begin{equation}\label{eq:means}
 N_j=\sum_iw_{ji},\quad a_{ji}=w_{ji}/N_j,\quad
 \mu_j=\sum_i a_{ji}x_i,\quad \bar y_j=\sum_i a_{ji}y_i,\quad
 n_{{\rm eff},j}=\frac{1}{\sum_i a_{ji}^2}.
\end{equation}
For $q_{jsi}=\phi_{js}^\top(x_i-\mu_j)$, the numerical implementation
also subtracts the residual weighted mean
$\bar q_{js}=\sum_i a_{ji}q_{jsi}$. The normalized statistics are
\begin{align}
 I_{js}&=\sum_i a_{ji}(q_{jsi}-\bar q_{js})(y_i-\bar y_j),
 \label{eq:I}\\
 U_{js,:}&=\sum_i a_{ji}(q_{jsi}-\bar q_{js})(x_i-\mu_j)^\top.
 \label{eq:U}
\end{align}
Exact arithmetic gives $\bar q_{js}=0$. Explicit recentering improves
stability when features have a large offset. The response mean stored as
\code{statistics.S} is exactly $\bar y_j$; it is also used in outer-step
selection.

Define empirical weighted covariance and cross-covariance
\begin{equation}\label{eq:cov}
 \Sigma_j=\sum_i a_{ji}(x_i-\mu_j)(x_i-\mu_j)^\top,\qquad
 c_j^{xy}=\sum_i a_{ji}(x_i-\mu_j)(y_i-\bar y_j).
\end{equation}
Then the statistics have the concise interpretation
\begin{equation}\label{eq:sketch-cov}
 U_j=\Phi_j\Sigma_j,\qquad I_j=\Phi_j c_j^{xy}.
\end{equation}
The code computes these products directly from observations and directions;
it need not form $\Sigma_j$.

\begin{proposition}[Why a common gradient subspace gives moment equations]
If the responses within a neighborhood obey the affine model
$y_i=b_j+(Yg_j)^\top x_i+e_i$, then
\begin{equation}\label{eq:moment-model}
 I_j=U_jYg_j+\Phi_j\sum_i a_{ji}(x_i-\mu_j)e_i.
\end{equation}
In particular, exact affine responses with gradient in $\col(Y)$ satisfy
$I_j=U_jYg_j$ exactly.
\end{proposition}
\begin{proof}
Subtract the weighted response mean. The intercept disappears and
$c_j^{xy}=\Sigma_jYg_j+\sum_i a_{ji}(x_i-\mu_j)e_i$.
Multiply by $\Phi_j$ and use \eqref{eq:sketch-cov}.
\end{proof}
For smooth $f$, a local Taylor expansion motivates this affine model;
its remainder and observational noise contribute to $e_i$. This motivation
does not remove finite-bandwidth bias.

For isotropic directions and a fixed vector $v$, independent of the sketch,
$\E\norm{\Phi_j v}^2=(p/d)\norm{v}^2$. However, the algorithm uses a
\emph{realized finite sketch} and optimizes a basis against it. Taking an
expectation does not make its optimizer equal to the optimizer of a full
covariance criterion. Even a full covariance-moment residual is not the
same objective as observation-level local linear least squares.

\paragraph{Legacy normalization.}
The explicit \code{legacy} path returns $N_jI_j,N_jU_j$ and passes no mass
to the solver, which then uses unit weights. At zero solver ridge this
weights the normalized residual by $N_j^2$, whereas the new path below
uses $N_j$. Legacy is a different normalization, not an algebraically
identical implementation of the new estimator.

\section{The frozen profiled objective on the Grassmann manifold}
\label{sec:profile}
For one outer step, freeze all $U_j,I_j,N_j$. The Grassmann solver targets
\begin{equation}\label{eq:profile}
 \boxed{\quad F_\tau(Y)=\sum_j N_j\min_{g\in\R^m}
 \left\{\norm{I_j-U_jYg}^2+\tau\norm{g}^2\right\},\quad
 Y^\top Y=I_m.\quad}
\end{equation}
Here $\tau=\code{grassman.solve.local_ridge}$ is zero by default.
It is \emph{separate} from $\lambda_{\rm init}$ in
\eqref{eq:init-fit}. The adapter does not pass the configuration's
\code{local_ridge} into this solver. To set a positive $\tau$, pass it
explicitly to \code{ADP_solver(solve, local_ridge=...)}.

For $M_j=U_jY$ and $H_j=M_j^\top M_j+\tau I_m$, the local solution is
\begin{equation}\label{eq:local-solution}
 g_j=M_j^\dagger I_j\quad(\tau=0),\qquad
 g_j=H_j^{-1}M_j^\top I_j\quad(\tau>0),\qquad
 r_j=I_j-M_jg_j.
\end{equation}
The first expression means the minimum-norm least-squares solution.
Inverse notation in the second expression specifies a mathematical map;
the code uses QR/SVD factors rather than forming $H_j^{-1}$ or normal
equations. Positive $\tau$ is a statistical ridge modification of the
profile, whereas a damping parameter used to generate a numerical step
is not.

For an orthogonal $Q\in\R^{m\times m}$, replacing $Y$ by $YQ$ and
$g_j$ by $Q^\top g_j$ leaves \eqref{eq:profile} unchanged. Consequently,
the natural parameter is $\col(Y)\in\Gr(m,d)$, with representative $Y$.
This objective profiles all local coefficients at each candidate basis;
it is distinct from a rank-constrained matrix problem with fixed $g_j$.

\subsection{Local linear algebra and rank boundaries}
Local ridge fits are implemented by augmenting the design:
\begin{equation}
 \widetilde M_j=\begin{bmatrix}M_j\\\sqrt\tau I_m\end{bmatrix},\qquad
 \widetilde I_j=\begin{bmatrix}I_j\\0\end{bmatrix}.
\end{equation}
The SVD path retains singular values larger than
$\epsilon_{\rm mach}\max(p,m)\sigma_{j,1}$. This is a numerical-rank rule;
when $\tau>0$, the factor in this cutoff still uses $p,m$, not $p+m,m$.
A fast reorthogonalized QR path is used for safe $m\le2$ systems, with
fallback to SVD for sensitive batches. The smooth flag additionally
requires full local rank, separation from severe conditioning, and an
absolute singular-value guard. Floating-point rank truncation means that
the formal exact profile and its computed realization can differ near
a rank threshold.

At $\tau=0$, local rank can change with $Y$. Derivatives that assume full
column rank are not used there. A small computed gradient at a rank boundary
is not reported as a smooth stationarity certificate.

\subsection{Horizontal gradient}
On the smooth region, the envelope theorem gives
\begin{equation}\label{eq:gradient}
 \nabla_YF_\tau=-2\sum_jN_jU_j^\top r_jg_j^\top,
 \qquad G=(I_d-YY^\top)\nabla_YF_\tau.
\end{equation}
To verify the cancellation, differentiate the local objective:
the coefficient variation contributes
$2(\tau g_j-M_j^\top r_j)^\top dg_j=0$ by the local normal condition.
Only the direct variation through $M_j$ remains. The horizontal tangent
space is $\{Z:Y^\top Z=0\}$; projecting gives $G$.
The implementation applies the projection as $G_E-Y(Y^\top G_E)$,
without forming $I_d-YY^\top$. Outside the smooth region it uses the same
algebraic expression as a proposal direction and withholds the certificate.

The stopping normalization is
\begin{equation}\label{eq:certificate}
 s_I=\max\left(1,\sum_jN_j\norm{I_j}^2\right),\qquad
 \gamma=\norm{G}_F/s_I.
\end{equation}
The default tolerance is $10^{-6}$. It is this normalized horizontal
gradient, together with applicability of smooth stationarity, that controls
the inner \code{converged} flag.

\section{The current inner update: a reduced Gauss--Newton core}
\label{sec:gn}
\subsection{Choosing the search space}
Compute a thin SVD $G=V_G\diag(\sigma)A^\top$. Let $r_G$ be its
numerical rank, capped by $\min(m,d-m)$. Unless a fixed \code{rank} is
provided, choose the smallest $q\le r_G$ such that
\begin{equation}\label{eq:energy}
 \sum_{a=1}^q\sigma_a^2\ge(1-\eta)\sum_a\sigma_a^2,
 \qquad\eta=\code{energy_tol}=0.05,
\end{equation}
capped at the available rank. With a fixed rank, use its minimum with
$r_G$. Reproject and QR-orthonormalize the first $q$ left vectors to obtain
$V\in\R^{d\times q}$, $Y^\top V=0$. Recompute the actual reduced gradient
$V^\top G$ after QR; its signs need not match the original SVD signs.

The search directions are $Z=VK$, with $K\in\R^{q\times m}$. All
$qm$ entries of $K$ are optimized together. This is a full core within
the selected left search space, not the full Grassmann tangent space,
whose dimension is $m(d-m)$. The parameter \code{rank} controls this
search restriction, not the estimated EDR dimension $m$.

Cache $M_j=U_jY$ and $W_j=U_jV$ for the current step only. These operator
actions allow all trial bases to be evaluated without rereading the
ambient dimension in each profile trial.

\subsection{Differentiating the profiled residual, including coefficient refits}
For a tangent perturbation $Z=VK$, write $T_j=U_jZ=W_jK$.
Differentiating $H_jg_j=M_j^\top I_j$ yields
\begin{align}
 H_j\,dg_j&=T_j^\top r_j-M_j^\top T_jg_j,\label{eq:dg}\\
 dr_j&=-T_jg_j-M_j\,dg_j.\label{eq:dr}
\end{align}
The term $T_j^\top r_j$ is essential: holding coefficients fixed or
discarding that term gives a different Jacobian.

Stack the weighted augmented residuals
\begin{equation}\label{eq:aug-residual}
 R(Y)=\left(\sqrt{N_j}
     \begin{bmatrix}r_j\\-\sqrt\tau g_j\end{bmatrix}\right)_{j=1}^J,
 \qquad \norm{R(Y)}^2=F_\tau(Y),
\end{equation}
omitting the bottom rows when $\tau=0$. Equations \eqref{eq:dg}--\eqref{eq:dr}
give the Jacobian $\mathcal J$ with respect to the row-major vector
$\operatorname{vec}(K)$. In particular,
\begin{equation}\label{eq:jac-gradient}
 2\mathcal J^\top R=\operatorname{vec}(V^\top G).
\end{equation}
The code constructs this small-column Jacobian, including its coefficient
cross terms; it does not build a Jacobian for all $m(d-m)$ tangent variables.

\subsection{Damped least squares and the polar retraction}
The numerical proposal solves
\begin{equation}\label{eq:gn-ls}
 \min_K\norm{\mathcal J\operatorname{vec}(K)+R}^2
       +\lambda_d\norm{K}_F^2
 =\min_K\left\lVert
 \begin{bmatrix}\mathcal J\\\sqrt{\lambda_d}I_{qm}\end{bmatrix}
 \operatorname{vec}(K)-\begin{bmatrix}-R\\0\end{bmatrix}
 \right\rVert^2.
\end{equation}
The equality denotes two forms of the same minimization, up to the same
constant and with identical objective values as written.
The augmented least-squares solve avoids forming
$\mathcal J^\top\mathcal J$. On the smooth full-center path,
$\lambda_d=\code{lambda_prox}=0.05$ initially. This parameter arrives from
\code{ADP_Config.lambda_penalty}; despite that configuration name, it is
\emph{only damping} for this solver, not a permanent penalty on
$Y-Y_{\rm old}$ or on projectors in \eqref{eq:profile}.

The spectral norm of $K$ is capped at $\tan(\theta_{\max})$, where
$\theta_{\max}=0.5$ radians. A polar retraction gives
\begin{equation}\label{eq:polar}
 Y_+(t)=(Y+tVK)(I_m+t^2K^\top K)^{-1/2},\qquad
 M_{j,+}(t)=(M_j+tW_jK)(I_m+t^2K^\top K)^{-1/2}.
\end{equation}
The inverse square root is computed in dimension $m$.
Since $Y^\top V=0$ and $V^\top V=I_q$,
$Y_+(t)^\top Y_+(t)=I_m$. The principal rotation angles are
$\arctan(t\sigma_a(K))$, so the cap controls the largest angle at $t=1$.

The proposed direction must have negative actual slope
$s=\langle V^\top G,K\rangle_F$. Trial lengths
$t=1,1/2,\ldots,2^{-15}$ are accepted only if the recomputed profile obeys
\begin{equation}\label{eq:armijo}
 F_\tau(Y_+(t))\le F_\tau(Y)+10^{-4}t s.
\end{equation}
Local coefficients are reoptimized at every trial. If no length is accepted,
damping is multiplied by ten and another least-squares proposal is tried,
up to eight damping attempts. A failed proposal triggers a spectral fallback.
This is a bounded descent procedure, not an exact minimization of the
Grassmann objective.

\section{Rank loss, constant centers, and descent fallbacks}\label{sec:guards}
\subsection{Exactly zero operators}
If $U_j=0$ exactly, the minimum-norm coefficient is $g_j=0$ and that
center contributes the known constant $N_j\norm{I_j}^2$ to
\eqref{eq:profile}, for any $\tau\ge0$. With
\code{constant_centers=True}, a mixture of zero and nonzero operators is
reduced to the active centers for the solve. The code restores zero
coefficient rows, constant objective/history terms, and rank counts.
It rescales the active tolerance by $s_I/s_{I,\rm active}$, preserving
the original stopping threshold in \eqref{eq:certificate}. This is exact
zero testing, not pruning small operators. An entirely zero batch follows
the direct path and is not given a false smooth certificate at $\tau=0$.

\subsection{Relative Gauss--Newton with intrinsic local rank loss}
At $\tau=0$, if some $M_j$ are rank deficient, the current
\code{relative_gn=True} branch checks whether their numerical ranks satisfy
\begin{equation}\label{eq:intrinsic}
 \rank(U_j)=\rank(U_jY)<m.
\end{equation}
In exact linear algebra this equality implies
$\col(U_jY)=\col(U_j)$, so the center's least-squares value is locally
constant as long as that equality persists. The residual is orthogonal to
$\col(U_j)$ and its ambient gradient contribution vanishes.

The implemented test uses SVD numerical ranks, not an exact symbolic
identity. It therefore justifies a \emph{proposal model}: if every deficient
center passes the rank test, and at least one remaining center exists with
a smooth profile, GN is formed using those regular centers only. The
initial damping is then scaled relatively:
\begin{equation}\label{eq:relative}
 \lambda_d=\code{lambda_prox}\frac{\norm{\mathcal J_{\rm reg}}_F^2}{qm}.
\end{equation}
The direction space and slope still use the full algebraic gradient.
Armijo acceptance and the final objective use \emph{all} centers.
Nonpositive or nonfinite relative damping invalidates the proposal.
No full-rank derivative is applied to a deficient center, and numerical
rank agreement is not promoted into a smooth full-profile certificate.

\subsection{Value-based rank-one rotations}
If an admissible GN model cannot be obtained, take a leading gradient
singular pair, with unit $a\in\R^m$ and $v\perp\col(Y)$ oriented so that
$s=v^\top Ga<0$. The rank-one Grassmann rotation is
\begin{equation}\label{eq:rank-one}
 Y(\theta)=Y+\bigl[Ya(\cos\theta-1)+v\sin\theta\bigr]a^\top.
\end{equation}
For $u_j=M_ja$ and $w_j=U_jv$, its cached action is
\begin{equation}
 M_j(\theta)=M_j+
 \bigl[u_j(\cos\theta-1)+w_j\sin\theta\bigr]a^\top.
\end{equation}
This curve preserves orthonormality. All trial values are evaluated with
fresh local coefficient fits, so a nonsmooth rank transition does not
require the full-rank Jacobian formula.

The default \code{quadratic_trials=3} starts at $\theta_{\max}$. From a
trial value it estimates
\begin{equation}\label{eq:quadratic}
 \widehat c=\frac{2[F_\tau(Y(\theta))-F_\tau(Y)-\theta s]}{\theta^2},
 \qquad
 \theta_{\rm prop}=\operatorname{clip}
 \left(-s/\widehat c,\theta_{\max}2^{-19},\theta_{\max}\right),
\end{equation}
when $\widehat c$ is positive and finite. Of up to three such trials, the
best one actually satisfying Armijo is retained. If none qualifies, the
code tries up to twenty dyadic angles. Curvature interpolation is a search
heuristic; the true profile decides acceptance.

If both \code{relative_gn=False} and \code{quadratic_trials=0}, a retained
compatibility branch uses a scalar Schur-complement angle profile, followed
by full-refit Armijo checks. It is not the current default path.

\subsection{Spectral fallback and workspace guard}
If the GN workspace estimate exceeds \code{workspace_bytes}, or GN cannot
produce an accepted step, a spectral descent curve is used. Let
$V^\top G=L\diag(\sigma)A^\top$ and $\omega_a=\sigma_a/\sigma_1$.
With $\widehat V=VL$, it is
\begin{equation}\label{eq:spectral}
 Y(t)=Y+YA\diag(\cos(t\omega)-1)A^\top
          -\widehat V\diag(\sin(t\omega))A^\top.
\end{equation}
Here $A$ has $q$ orthonormal columns. The derivative has negative slope
$-\sum_a\sigma_a\omega_a$, and up to twenty dyadic lengths starting at
$\theta_{\max}$ are checked using the same true-profile Armijo rule.
Workspace controls the GN scratch space, not the full input $U$ or total
process memory.

\section{Returning from the inner solve and updating the outer spectrum}
\label{sec:canonical}
The default inner budget is at most twenty accepted iterations, with no
separate smooth-region cap. The solver can stop earlier with
\code{stationary}, \code{rank_boundary_stationary}, or
\code{line_search_failed}; exhausting the budget gives \code{max_steps}.
Before returning, it recomputes $U_jY$, the profile, and the gradient from
the live operators rather than trusting accumulated cached trial actions.

The result contains the row basis $P=Y^\top$, local coefficients
$g_j^\top$ as rows, and diagnostics. Inner \code{converged=True} requires
the normalized gradient threshold and applicability of the smooth
certificate. Reaching a finite step cap is not convergence.

For the normalized outer algorithm, assemble $C\in\R^{J\times m}$
with rows $\sqrt{N_j}g_j^\top$, and compute
\begin{equation}\label{eq:canonical}
 C=L_C\diag(s_1,\ldots,s_m)R_C^\top,\qquad
 B_{k+1}=YR_C,\qquad
 \ell_{k+1,a}=s_a^2/s_1^2.
\end{equation}
Column signs are made deterministic by making each column's largest
absolute entry positive. The coefficients are rotated by the inverse
basis rotation, including the signs, to preserve every $Yg_j$.
The required coefficient rank is checked; if it is below $m$, the update
is rejected as non-identifiable rather than silently reducing $m$.

This is the eigendecomposition of the ambient gradient surrogate
\begin{equation}
 \sum_jN_j(Yg_j)(Yg_j)^\top
       =Y(C^\top C)Y^\top
\end{equation}
using only the small factor $C$. It changes the representative basis and
localization spectrum, while preserving the estimated subspace at this
point. The spectrum then influences the \emph{next} outer neighborhoods.
For the legacy path the coefficient SVD is unweighted.

\section{Outer-step scoring, stopping, and complete pseudocode}
\label{sec:outer}
\subsection{What the selection score actually measures}
At outer step $k$ the code records
\begin{equation}\label{eq:selection}
 E_k=\sum_{j=1}^J\bigl(y_{i_j}-\bar y_j^{(k)}\bigr)^2,
\end{equation}
where $\bar y_j^{(k)}$ was computed from the neighborhoods determined by
the \emph{incoming} $B_k,\ell_k,h_k,\alpha_k$. After the solver update,
the record associates this score with the newly computed $B_{k+1}$.
There is no second statistics pass to score the updated basis at the same
scale. Thus $E_k$ is neither the solver profile nor a fitted local linear
prediction using $g_j$.

The default \code{select_step="best"} returns the updated basis attached
to the smallest recorded $E_k$, keeping the first in a tie; \code{last}
returns the final update. With the default training protocol this is a
same-sample smoothing criterion. With excluded center indices it is a
center-response holdout criterion, but repeated selection across scales
still does not constitute untouched scientific validation.

The outer loop stops on the explicit \code{outer_steps} cap, when
$h_k/a<h_{\min}$, or when no feasible $\alpha$ can maintain the target
mass at the next scale. These stops are independent of inner stationarity.
The first step is executed even if $h_0=h_{\min}$. The initial basis alone
is not a candidate in the best-step selection.

\subsection{Complete outer algorithm}
The following uses zero-based $k$ and the normalized default path. Optional
training splits and displaced centers are applied before the first line
that constructs training distances.
\begin{enumerate}[label=\arabic*.,leftmargin=*]
\item Validate $X,y,m$ and configuration; resolve $J,p,N_{\rm lin}$.
Create random streams, select centers and training observations.
Compute training-to-center squared distances.
\item Estimate $B_0$ by local gradient PCA
(Section~\ref{sec:init}); set $\ell_0=\bm1$, $\alpha_0=1$.
Find $h_0\ge h_{\min}$ satisfying \eqref{eq:h-search}.
\item For $k=0,1,\ldots$:
\begin{enumerate}[label=(\alph*),leftmargin=*]
\item Draw or reuse $\Phi_j$. Force the orthogonal tensor at $k=0$;
otherwise use the configured tensor.
\item Build weights, masses, means, and normalized $I_j,U_j$ using
\eqref{eq:orthogonal}--\eqref{eq:U}. Keep $\bar y_j$ for scoring.
\item Call the inner Grassmann solver with initial row basis $B_k^\top$,
operators $U_j$, targets $I_j$, masses $N_j$, and configured damping.
\item Validate the returned basis and coefficients. Canonicalize by
\eqref{eq:canonical}, obtaining $B_{k+1},\ell_{k+1}$ and rotated coefficients.
\item Compute \eqref{eq:selection} from the already constructed means;
store the updated basis, spectrum, coefficients, score, geometry, and
inner diagnostics. Update the best record if its score is strictly lower.
\item Stop if the outer budget is reached. Put $h'=h_k/a$; stop if
$h'<h_{\min}$. Find $\alpha'$ by \eqref{eq:alpha-search} using
$B_{k+1},\ell_{k+1},h'$; stop if infeasible. Set
$h_{k+1}=h'$, $\alpha_{k+1}=\alpha'$.
\end{enumerate}
\item Return the best recorded update or the last update as configured.
Publish the column basis and project new observations with $XB_{\rm out}$.
\end{enumerate}

\subsection{Complete inner algorithm}
\begin{enumerate}[label=\arabic*.,leftmargin=*]
\item Validate shapes, finite data, positive masses, and solver parameters.
If enabled, eliminate exactly zero operators and remember their constants.
Orthonormalize the initial row basis and transpose it to $Y$.
\item Compute the local profile $M_j=U_jY$, $g_j,r_j,F_\tau$ by QR/SVD.
For up to the inner step budget:
\begin{enumerate}[label=(\alph*),leftmargin=*]
\item Compute the full algebraic horizontal gradient. If its normalized
norm is below tolerance, stop, distinguishing smooth stationarity from a
rank boundary.
\item Select $V$ and $q$ by the gradient SVD. If the full profile is
smooth, use all centers for GN; otherwise attempt the eligible intrinsic
rank model, or choose a rank-one value search.
\item Cache $U_jV$. If GN is admissible and fits the workspace, construct
\eqref{eq:dg}--\eqref{eq:gn-ls}, cap the tangent length, and try damped
polar steps with full-profile Armijo acceptance.
\item Use the rank-one or spectral fallback when appropriate. If no trial
is accepted, stop with \code{line_search_failed}; otherwise replace $Y$
and its profile by the accepted state and append the objective history.
\end{enumerate}
\item Recompute the final profile and gradient from live $U_jY$. Set the
certificate and diagnostics, restore constant centers, and return $Y^\top$.
\end{enumerate}

\section{Numerical interpretation, resource bounds, and limitations}
\label{sec:limits}
\subsection{Three distinct quantities and their diagnostics}
The solver reports
\begin{align}
 \code{profile_objective}&=F_\tau(Y),\\
 \code{loss}&=\tfrac12\sum_jN_j\norm{r_j}^2,\\
 \code{riemannian_gradient}&=\norm{G}_F/s_I.
\end{align}
The \code{profile_history} stores $F_\tau$ along accepted steps;
\code{loss_history} is half that history. At $\tau>0$,
\code{loss_history} includes half the ridge term, whereas the final
\code{loss} excludes it. At $\tau=0$ they have the expected common scaling.
The outer score $E_k$ in \eqref{eq:selection} is a third, distinct criterion.

Other diagnostics include accepted steps, attempted step ranks, GN solves,
profile evaluations, rank-guard and workspace fallbacks, intrinsic GN
proposals, final rank loss, orthogonality error, stop reason, and
\code{stationarity_applicable}. The reported method name can be
\code{core_gn_relative} and the angle backend \code{refit_quadratic}
under current defaults; diagnostic labels do not select a separate
public algorithm.

\subsection{Actual storage and numerical execution}
The main persistent sketch arrays cost $O(Jpd)$ for $U$ and directions,
plus $O(Jp)$ for $I$. The shared outer loop also precomputes a
$J\times n_{\rm tr}$ distance array, so its memory is not purely sparse.
Weights and their observation projections are processed in center blocks;
compact neighborhoods can use exact nonzero-neighbor statistics. Neither
a $J\times n_{\rm tr}\times d$ differences tensor nor a dense ambient
projector is required.

Inside the solver, caches cost $O(Jp(m+q))$. The GN Jacobian has
$J(p+m\mathbf1_{\tau>0})$ rows and $qm$ columns. Its dominant work is
small-column least squares, in addition to operator actions of order
$O(Jpd(m+q))$ per refreshed search space. The default GN workspace cap is
$16$ MiB and the ambient-gradient scratch is chunked at approximately
$4$ MiB; these limits do not cap total memory or input storage.
Arithmetic is float64. NumPy and CuPy follow the same mathematical
objective and acceptance logic; the GPU path additionally QR-compresses
some local SVD and augmented GN designs. Returned basis and coefficient
arrays are on the host. No performance ordering between CPU and GPU
follows from these formulas.

\subsection{What follows mathematically, and what does not}
\begin{itemize}[leftmargin=*]
\item Weighted centering removes local intercepts. Equations
\eqref{eq:sketch-cov}, \eqref{eq:moment-model}, and gauge invariance are
exact arithmetic identities.
\item On a smooth full-rank profile, the gradient, coefficient derivative,
Jacobian relation, and polar orthogonality formulas are exact. Numerical
rank truncation and floating-point errors require the stated safeguards.
\item Accepted inner steps satisfy the implemented sufficient-decrease
test for their frozen profile. A finite twenty-step budget does not imply
stationarity; smooth stationarity does not imply a global minimum.
\item Intrinsic-rank tests, selected gradient search spaces, relative
damping, and bounded quadratic trials are optimizer choices. Their
acceptance is checked on the full profile, but they are not a convergence
proof at rank boundaries.
\item Noise, local Taylor remainder, finite direction sketches, and
adaptive neighborhoods all affect statistical recovery. Neither low
profile loss nor small center-mean SSE establishes consistency or a
finite-sample recovery rate.
\end{itemize}
If truth is available in an experiment, use basis-invariant metrics, for
example $\norm{\widehat B\widehat B^\top-B_\star B_\star^\top}_F$
or principal angles, and report them separately from time, memory, and
solver certificates. Ground-truth quality callbacks in the shared engine
record diagnostics; they do not determine the fitted update or score.

\appendix
\section{Current solver defaults and parameter ownership}\label{sec:defaults}
\begin{center}
\begin{tabular}{ll}
\toprule
\code{grassman.solve} argument & Default\\
\midrule
\code{max_steps}, \code{tol} & $20$, $10^{-6}$\\
\code{lambda_prox}, \code{local_ridge} & $0.05$, $0$\\
\code{rank}, \code{energy_tol} & automatic, $0.05$\\
\code{max_angle}, \code{workspace_bytes} & $0.5$ radians, $16$ MiB\\
\code{relative_gn}, \code{quadratic_trials} & \code{True}, $3$\\
\code{constant_centers}, \code{smooth_max_steps} & \code{True}, no separate cap\\
\bottomrule
\end{tabular}
\end{center}
An explicit \code{smooth_max_steps} can cap an initially smooth solve
when relative GN or quadratic trials are enabled; it is not active by
default. $\code{lambda_prox}>0$ is required even though
\code{ADP_Config.lambda_penalty} permits zero for other solvers.
Do not pass \code{lambda_prox} or \code{mass} as fixed
\code{ADP_solver} settings for this facade: they are already passed by
the model and duplicates are rejected.

The production \code{solve} exposes no public \code{method} or
\code{angle_backend} selector. Alternative flat/full tangent GN methods
belong to \code{grassman_exp.py}; \code{block_rank_k} is a separate
experimental entry point. They are not part of this document's main
algorithm. The explicit older control
\begin{lstlisting}[language=Python]
ADP_solver(solve, max_steps=5, relative_gn=False,
           quadratic_trials=0, constant_centers=False)
\end{lstlisting}
changes optimization settings while retaining the same frozen objective
at the same solver ridge.

\section{Implementation source map and reproducibility}\label{sec:sources}
This is a description of the live source snapshot, not a reinterpretation
of older manuscript formulas. The following paths and named functions
identify the authoritative implementation components.
\begin{longtable}{@{}>{\raggedright\arraybackslash}p{0.48\textwidth}
  >{\raggedright\arraybackslash}p{0.48\textwidth}@{}}
\toprule
Source & Responsibility\\
\midrule
\endhead
\path{ADP/core/multi/ADP_multi_index.py} &
Facade, explicit solver integration, public basis and transform.\\[3pt]
\path{ADP/core/ADP_Config.py} &
Outer defaults, kernel, configuration validation.\\[3pt]
\path{ADP/core/ADP_Solver.py}, \code{fit_current} &
Mass normalization and damping forwarding; separation of solver settings.\\[3pt]
\path{ADP/engine/common/index_fit.py},
\code{validate_sizes}, \code{_initial_index},
\code{canonical_index}, \code{fit_index} &
Sizes, initialization, full outer loop, spectrum, score chronology, stops,
best/last selection, coefficient rotation.\\[3pt]
\path{ADP/engine/common/initialize.py},
\code{initialize_basis_local_with_spectrum} &
Local ridge gradients, weighted gradient PCA and alternative initializers.\\[3pt]
\path{ADP/engine/common/calculus.py},
\code{search_bandwidth}, \code{calculate_alpha_k},
\code{generate_multi_proj}, \code{generate_isotropic_proj} &
Bandwidth, anisotropy, compact-support screening and sketch laws.\\[3pt]
\path{ADP/engine/common/weights.py}, \code{calculate_multi_weight} &
Principal/residual geometry, blockwise kernel evaluation.\\[3pt]
\path{ADP/engine/common/statistic.py},
\code{calculate_statistics}, \code{_dense_moments}, \code{_local_block} &
Centered statistics, normalized/legacy scaling, response means, effective
sample sizes.\\[3pt]
\path{ADP/solver/grassman.py}, \code{_profile},
\code{_gradient}, \code{_core_jacobian}, \code{_polar},
\code{_intrinsic_model}, \code{solve} &
Complete current production inner solver.\\[3pt]
\path{tests/test_grassman.py},
\path{tests/test_grassman_combined.py} &
Independent local fits and finite differences, gauge/orthogonality,
rank boundaries, combined defaults and constant-center restoration.\\
\bottomrule
\end{longtable}
Build from the repository root with
\begin{lstlisting}
pdflatex -interaction=nonstopmode -halt-on-error \
  -output-directory=/tmp/adp-grassmann-tex tex/ADP_Grassmann_current.tex
# Repeat once for the table of contents and cross-references.
\end{lstlisting}
Create the output directory first. The adjacent verification record
\path{tex/ADP_Grassmann_current.verification.md} records the actual commands,
formula checks, focused tests, source hashes, and build outcome for this
snapshot. The accompanying PDF is generated from this TeX source.
\end{document}
